% =====================================================================
% SECÇÃO 5 — RESULTS AND DISCUSSION
% Reescrita com os dados REAIS da experiência de 16B (27 projetos, 486 módulos).
%
% NÚMEROS CONFIRMADOS (fonte: coverage_measured.csv, mutmut.csv,
% test_quality.csv, run_results/*.json — todos em results/deepseek-coder-v2_16b/)
%
% MARCADORES:
%   [[MUTMUT-ANSIBLE]] -> substituir quando o mutmut por-módulo terminar
%   [[QWEN]]           -> preencher com a run do Qwen2.5-Coder-32B
%   [[WILCOXON]]       -> preencher com p-values e tamanho de efeito
% =====================================================================

\section{Results and Discussion}

In this section we compare MARTA against the single-prompt baseline (Test4Py) and
the search-based generator (Pynguin) across the 27 projects (486 modules) of our
corpus. As described in Section~\ref{sec:implementation}, MARTA and the
single-prompt baseline share the identical Phase~1 semantic context, the same
underlying model, and a repair budget of $N=3$; the results below report the full
pipeline with its outer coverage-driven loop enabled (three generation rounds).

\subsection{Test Suite Yield and Oracle Strength}

A recurring claim in LLM-based test generation is that richer context yields more
executable tests. Our data show that, at this model scale, the decisive difference
between architectures lies not in \emph{how many} tests survive validation, but in
\emph{what those tests actually verify}.

Both LLM-based pipelines converge on suites of comparable size: MARTA retains
11{,}255 test functions and the single-prompt baseline 11{,}235, while Pynguin
produces 4{,}100. The distinguishing factor is oracle strength. As shown in
Table~\ref{tab:oracles}, only 0.8\% of MARTA's tests contain no assertion at all,
against 7.1\% for the single-prompt baseline and 38.3\% for Pynguin. The baseline's
rate is a direct, measurable consequence of its \emph{heuristic reduction} step,
which minimizes a test until it executes cleanly and thereby removes its assertions;
MARTA, by contrast, discards unrepairable tests outright rather than hollowing them
out (Section~3.2.4). Pynguin's much higher rate reflects a different mechanism: as a
search-based generator optimizing structural coverage, a substantial portion of its
test cases exist purely to execute code paths, with regression assertions attached
only where its assertion generator can produce them.

The same asymmetry appears in assertion quality. We classify an assertion as
\emph{structurally weak} when it cannot discriminate program behaviour: assertions
over constants, \texttt{isinstance}/\texttt{hasattr} checks, comparisons against
\texttt{type(x)} (including the \texttt{f}-string type signatures characteristic of
Pynguin), and assertions over mock internals such as \texttt{call\_count}. Under this
deliberately conservative criterion, which does \emph{not} penalize comparisons
against literal values, 13.1\% of MARTA's assertions are weak, against 17.0\% for the
baseline and 30.4\% for Pynguin.

\begin{table}[htpb]
\centering
\caption{Oracle strength of the generated suites (16B model, 27 projects).
``No assertion'' is the share of test functions containing no assertion or
\texttt{pytest.raises}; ``weak'' is the share of assertions that cannot
discriminate behaviour. Parametrized tests are expanded to their number of cases.}
\label{tab:oracles}
\begin{tabular}{lrrrr}
\hline
\textbf{Tool} & \textbf{Tests} & \textbf{Asrt/test} & \textbf{No assertion} & \textbf{Weak} \\
\hline
MARTA            & 11{,}255 & 1.85 & \textbf{0.8\%}  & \textbf{13.1\%} \\
Test4Py          & 11{,}235 & 1.96 & 7.1\%           & 17.0\% \\
Pynguin          &  4{,}100 & 2.84 & 38.3\%          & 30.4\% \\
\hline
\end{tabular}
\end{table}

\subsection{Structural Coverage}

Because the benchmark projects differ by two orders of magnitude in size (from a
single module in \texttt{codetiming} to 237 in \texttt{ansible}), the choice of
aggregation materially changes the ranking. We therefore report three complementary
aggregations rather than selecting one: the mean over projects, a module-weighted
mean (the unit of evaluation used in the CodaMosa benchmark, where each of the 486
modules is a benchmark instance), and a pooled measure computed as the ratio of all
covered lines and branches to all coverable ones, which is the ``overall'' coverage
reported by CoverUp. The pooled measure is the strictest of the three: it uses an
identical denominator (68{,}248 coverable elements) for all tools and involves no
weighting choice.

\begin{table}[htpb]
\centering
\caption{Line+branch coverage over the target modules under three aggregations
(16B model, 27 projects). Pooled coverage uses an identical denominator for all
three tools.}
\label{tab:coverage}
\begin{tabular}{lrrr}
\hline
\textbf{Aggregation} & \textbf{MARTA} & \textbf{Test4Py} & \textbf{Pynguin} \\
\hline
Pooled (overall)        & \textbf{40.6\%} & 37.1\% & 35.3\% \\
Module-weighted         & \textbf{41.8\%} & 38.0\% & 38.7\% \\
Mean over projects      & 49.1\%          & 44.9\% & \textbf{54.7\%} \\
Median over projects    & 42.1\%          & 43.3\% & \textbf{49.1\%} \\
\hline
\end{tabular}
\end{table}

MARTA leads under both the pooled and the module-weighted aggregation, and improves
over the single-prompt baseline under all four, winning 19 of the 27 pairwise
per-project comparisons. Under the mean over projects, however, Pynguin ranks first.
This inversion is informative rather than incidental. Averaging over projects assigns
equal weight to \texttt{codetiming} (one module) and to \texttt{ansible} (237
modules), and Pynguin's advantage is concentrated precisely in small, structurally
simple modules where an evolutionary search can exhaust the input space: it reaches
100\% on \texttt{sty} and 96.2\% on \texttt{codetiming}. On the largest and most
structurally complex subjects the ordering reverses: on \texttt{ansible}, which alone
accounts for 49\% of the benchmark's modules, MARTA attains 40.0\% against Pynguin's
27.7\%; on \texttt{flutils} the gap is 66.6\% against 33.6\%, and on \texttt{black}
Pynguin covers none of the target modules. This is the coverage plateau that
motivated LLM-augmented search in the first place \cite{lemieux2023codamosa}: search
saturates where the state space is large and object construction is non-trivial,
which is exactly the regime where semantic context helps.

\subsection{Fault-Revealing Capability}

High coverage is a well-known vanity metric: executing a line does not imply
detecting a fault in it. We therefore assess the generated suites with mutation
analysis. Because the projects differ in size, mutants are attributed to the module
they belong to, and modules for which a tool produced no usable test contribute their
mutants as survived rather than being excluded, so that a tool is not rewarded for
generating fewer tests.

Over the [[MUTMUT-ANSIBLE: 23]] projects for which all three tools yield a mutation
score, MARTA achieves a mean score of 26.5\% against 25.0\% for the single-prompt
baseline and 21.9\% for Pynguin, and obtains the highest score on 13 projects
against 6 and 7 respectively. The comparison with Pynguin is the more informative
one: Pynguin attains higher per-project coverage yet detects fewer injected faults,
which is consistent with the oracle measurements in Table~\ref{tab:oracles}. A suite
in which 38.3\% of tests assert nothing accumulates structural coverage without
acquiring the ability to detect a regression.

% [[MUTMUT-ANSIBLE]] Substituir o parágrafo acima e a tabela quando a medição
% por-módulo terminar: passa a 27 projetos, incluindo ansible.
\begin{table}[htpb]
\centering
\caption{Mutation score (16B model). [[MUTMUT-ANSIBLE: atualizar para 27 projetos]]}
\label{tab:mutation}
\begin{tabular}{lrrr}
\hline
\textbf{Tool} & \textbf{Mean} & \textbf{Median} & \textbf{Best on} \\
\hline
MARTA    & \textbf{26.5\%} & \textbf{20.5\%} & \textbf{13} projects \\
Test4Py  & 25.0\%          & 19.8\%          & 6 projects \\
Pynguin  & 21.9\%          & 17.0\%          & 7 projects \\
\hline
\end{tabular}
\end{table}

\subsection{Computational Cost}

Decoupling generation into two agents plus a repair loop could reasonably be expected
to multiply inference cost. It does not. Over the 27 projects, MARTA consumed
115.5M tokens and 42.2 GPU-hours, against 201.6M tokens and 98.7 GPU-hours for the
single-prompt baseline: roughly half the cost for a suite of equivalent size and
stronger oracles. Two design decisions account for this. First, the Assertion Agent
emits one complete file per function in a single inference call rather than one call
per scenario, so the planning call is amortized across all scenarios of a function.
Second, the outer loop regenerates only functions with remaining uncovered lines,
and Phase~1 is restored from cache when the source is unchanged, so no LLM budget is
spent re-deriving context. Pynguin, which performs no inference, is excluded from
this comparison; its cost is CPU search time (300\,s per module).

\subsection{Effect of Model Capacity}

% [[QWEN]] Preencher quando a run do Qwen2.5-Coder-32B terminar.
% Hipótese a testar: se a vantagem MARTA-vs-baseline CRESCE de 16B para 32B,
% isso suporta a tese de que arquiteturas agênticas escalam com a capacidade do
% modelo. Com dois pontos só se pode afirmar uma DIREÇÃO, não uma curva.
[[QWEN: resultados com Qwen2.5-Coder-32B e comparação com o 16B.]]

\subsection{Statistical Significance}

% [[WILCOXON]] Preencher após correr os testes.
[[WILCOXON: Wilcoxon signed-rank sobre os pares por-projeto para cobertura,
mutation score e força dos oráculos, com tamanho de efeito (Vargha--Delaney
$\hat{A}_{12}$).]]
